\chapter[Proofs for Finite-Depth Adaptive Formation]{Proofs for Finite-Depth\\Adaptive Formation}
\label{app:finite-depth-formation-proofs}

This appendix proves the estimates used in
\cref{thm:finite-depth-adaptive-formation}. The carrier, normalized
experiment, and declared kinetics are those of
\cref{ch:finite-depth-formation}. Throughout, the depth $L\ge2$ is fixed.
Constants denoted by $c_L,C_L$ are positive and may change between
displays; they depend on that depth and the fixed model data, but not on
$\epsilon$, $A$, $T$, or the descending strengths unless this dependence
is displayed. All limits have $\epsilon>0$ and $T>0$. The results concern
the stated family at one base point. \status{HYPOTHESIS}

The orthogonal decomposition is
$\mathcal N\oplus\mathcal H\oplus\mathcal K$. On
$M_p=\mathcal N_p\oplus\mathcal H$, write $H_p$ for the Gaussian
Schur precision, $k_p=\dim M_p$, and $d_p=k_0-k_p$ for the number of
eliminated continuous coordinates. The normalized loop vector is $h$.
The entire assignment $S$ has values $j,k$ in the fixed finite set
$\mathcal J$; its value evolves, whereas its coordinate set is retained
at every level. Locally, $\pi_j(\vartheta)$ denotes its generative
prior mass, $\pi_\vartheta^p(x,j)$ the restricted optimizer density,
and $q_t^p(x,j)$ the actual joint recognition density, with respect to
Lebesgue measure times counting measure. The physical posterior
$\boldsymbol\Pi_\vartheta^p$ remains a different law. The independent
equilibrium root factor is omitted from non-root integrals and appended
unchanged. These are local notation conventions. \status{DEFINITION}

In the proofs, $D(\nu\Vert\mu)=\KL(\nu\Vert\mu)$ denotes relative
entropy. The same symbol applied to a density means its law with
the stated reference measure. \status{DEFINITION}

\section{Exact marginals, rates, and parameter derivatives}
\label{sec:rg-formation-exact-proof}

\lemmaheading{Weighted geometry and exact parent score}
{lem:formation-parent-score}
For the declared triangle-chain experiment, all marked marginal laws
and inherited proposal rows are normalized and compose exactly. On the
compact parameter interval $1\le\vartheta\le2.1$, the exact parent
potentials and complete model covectors have the forms established
below, including the eliminated Gaussian scalar and the source-prior
score. Their derivatives obey uniform energy-weighted bounds. These
statements are established for every fixed finite depth and all
sufficiently small positive $\epsilon$. \status{ESTABLISHED}

\emph{Proof.}
Let $B_s$ be the residual rows of level $s$. Independence follows
directly from their columns: a linear relation first vanishes on every
independent link column, forcing all internal-contrast and bridge
coefficients to vanish, and then vanishes on the model columns, forcing
all belief--model discrepancy coefficients to vanish. The normalized
loop rows are orthogonal to these rows. Choose fixed orthonormal flag
columns $\mathsf Q_s$ for
$\mathcal N_{s-1}\cap\mathcal N_s^\perp$. The matrix
$\mathsf B=B[\mathsf Q_1\ \cdots\ \mathsf Q_L]$ is block lower triangular with
invertible diagonal blocks, as is its inverse. If
$D_J=\operatorname{diag}(\sqrt{J_s}I)$, every nonzero block $(s,r)$ of
$D_J\mathsf B D_J^{-1}$ is multiplied by
$\sqrt{J_s/J_r}\le1$. Its Frobenius norm and the corresponding norm of
its inverse are therefore bounded by fixed graph-dependent constants.
Consequently
\begin{equation}
 c_L\sum_{s=1}^L J_s|w_s|^2
 \le w^TH_{\mathcal N}w
 \le C_L\sum_{s=1}^L J_s|w_s|^2,
 \qquad
 H|_{M_0}=H_{\mathcal N}\oplus J_1I_{\mathcal H}.
 \label{eq:rg-formation-weighted-proof}
\end{equation}
The Schur quadratic is the infimum of this quadratic over the
eliminated coordinates. Minimizing both bounds proves their retained
versions. In particular $H_p\succeq h_LI$ for a fixed $h_L>0$, since
$J_L=1$. The normal part of $H_p$ is bounded above by
$C_LJ_{p+1}I$, with $J_{L+1}=0$; the loop part still has weight $J_1$.

For a cumulative hidden/retained split $(z,x)$ of $M_0$, write
$F_p=H_{EE}$ and $G_p=H_{ER}$. Completing the square gives the
conditional-mean lift and its two identities
\begin{equation}
 L_px=(-F_p^{-1}G_px,x),\qquad
 L_p^THL_p=H_p,\qquad L_p^THz=0.
 \label{eq:rg-formation-gaussian-lift}
\end{equation}
Here and below $z$ in the last expression is embedded in the hidden
subspace. The centered hidden Gaussian has covariance
$(T/\vartheta)F_p^{-1}$, so its mean does not depend on $\vartheta$.
Let $\gamma_{p,\vartheta}$ denote this centered law. Exact integration
of the original source factor gives
\begin{align}
 V_{p,j}(x,\vartheta)
 &=\frac{\vartheta}{2}x^TH_px+C_p^G(\vartheta)
       +T\psi_{p,j}(x,\vartheta),
 \label{eq:rg-formation-parent-potential-proof}\\
 \psi_{p,j}(x,\vartheta)
 &=-\log\int e^{-\phi_j(L_px+z)}
                  \gamma_{p,\vartheta}(dz),
 \qquad 0\le\psi_{p,j}\le B_0,
 \label{eq:rg-formation-parent-tilt-proof}\\
 C_p^G(\vartheta)
 &=-\frac{Td_p}{2}\log\frac{2\pi T}{\vartheta}
       +\frac T2\log\det F_p,
 \qquad (C_p^G)'(\vartheta)=\frac{Td_p}{2\vartheta}.
 \label{eq:rg-formation-hidden-scalar-proof}
\end{align}
At $p=0$ the integral is absent and $C_0^G=0$. These formulas integrate
over whole Euclidean fibers. Positivity and Gaussian domination
justify iterated integration and differentiation on compact parameter
intervals. In particular the positive finite normalizer is retained,
and repeated integration agrees with direct integration.

The normalized non-root optimizer can also be written
\begin{equation}
 \pi_\vartheta^p(x,j)
 =Z_p(\vartheta)^{-1}e^{-\psi_{p,j}(x,\vartheta)}
       G_\vartheta^p(x)\pi_j(\vartheta),\qquad
 G_\vartheta^p=N\bigl(0,(T/\vartheta)H_p^{-1}\bigr),\qquad
 e^{-B_0}\le Z_p\le1.
 \label{eq:rg-formation-normalized-parent-proof}
\end{equation}
The microscopic proposal probability for the complete assignment is
\begin{equation}
 b_k(x,\vartheta)
 =\frac{\pi_k(\vartheta)e^{-\phi_k(x)}}
        {\sum_j\pi_j(\vartheta)e^{-\phi_j(x)}}.
 \label{eq:rg-formation-full-source-row-proof}
\end{equation}
It is the product of the normalized receiver rows at this fixed
configuration. The exact inherited row is
\begin{equation}
 R^p_{jk}(x,\vartheta)
 =\E_{\pi_\vartheta^0}
       [b_k(X,\vartheta)\mid P_pX=x,S=j].
 \label{eq:rg-formation-inherited-row-proof}
\end{equation}
Every row sums to one, including the self entry, and
$R^p_{jk}\ge e^{-B_0}\pi_k(\vartheta)>0$. The conservative
generator has diagonal $R^p_{jj}-1$, rather than $R^p_{jj}$. Its full
proposal clock has rate one and changes no continuous coordinate.
The microscopic identity
$\pi_\vartheta^0(X,j)b_k=
\pi_{\vartheta,X}^0(X)b_jb_k$ is symmetric in $j,k$.
Integrating over the same retained fiber proves
\begin{equation}
 \pi_\vartheta^p(x,j)R^p_{jk}(x,\vartheta)
 =\pi_\vartheta^p(x,k)R^p_{kj}(x,\vartheta).
 \label{eq:rg-formation-parent-balance-proof}
\end{equation}
The conditional-expectation tower proves composition of these rows
under another fixed retention. Separate integration of the physical,
optimizing, and actual joint laws through the same channel proves
their respective probability composition identities. No identification
of those laws is needed. In particular the retained actual source
conditional is not replaced by the optimizing conditional.

The fixed selected-source potential is
$\phi_j=\delta\sum_i f(\sqrt{J_1/T}r_{i,j_i})$,
where $f(t)=1-e^{-t^2/2}$. Thus $|f'|,|f''|\le1$.
Write $D_{\mathrm{int}}$ for the matrix with rows
$\sqrt{J_1}b_e$ on internal edges. Then
$D_{\mathrm{int}}^TD_{\mathrm{int}}\preceq H$. An edge is selected
$k_e\in\{0,1,2\}$ times and $\sum_e k_e=n$. Expressing the force as
$\delta\sqrt T D_{\mathrm{int}}^Tc$, with
$c_e=k_ef'(\sqrt{J_1/T}r_e)$, gives $|c|^2\le2n$. Hence, on $M_0$,
\begin{equation}
 |H^{-1/2}T\nabla\phi_j|\le B_*\sqrt T,
 \qquad
 -2\delta H\preceq T\nabla^2\phi_j\preceq2\delta H,
 \qquad B_*=\delta\sqrt{2n}.
 \label{eq:rg-formation-original-relative-bounds}
\end{equation}
For the marginal derivative, keep the centered Gaussian variable in
\cref{eq:rg-formation-parent-tilt-proof} fixed. In its normalized tilted
law $\nu_{p,x,j}$, differentiation gives
\begin{align}
 T\nabla\psi_{p,j}
 &=\E_{\nu_{p,x,j}}L_p^T(T\nabla\phi_j),
 \label{eq:rg-formation-marginal-force-proof}\\
 T\nabla^2\psi_{p,j}
 &=\E_{\nu_{p,x,j}}L_p^T(T\nabla^2\phi_j)L_p
       -T^{-1}\operatorname{Cov}_{\nu_{p,x,j}}
                    (L_p^TT\nabla\phi_j).
 \label{eq:rg-formation-marginal-hessian-proof}
\end{align}
The covariance has not been omitted. The isometry
$H^{1/2}L_pH_p^{-1/2}$ and
\cref{eq:rg-formation-original-relative-bounds} bound it by
$B_*^2H_p$. With $\theta_*=2\delta+2n\delta^2<1/2$, this proves
\begin{equation}
 |H_p^{-1/2}T\nabla\psi_{p,j}|\le B_*\sqrt T,
 \qquad
 (\vartheta-\theta_*)H_p
 \preceq\nabla^2V_{p,j}
 \preceq(\vartheta+2\delta)H_p.
 \label{eq:rg-formation-parent-relative-bounds}
\end{equation}
Evenness of the original source factor and centered Gaussian
integration also imply $\psi_{p,j}(-x,\vartheta)=
\psi_{p,j}(x,\vartheta)$. Every non-root sheet therefore has its
unique minimum at zero.

The hidden Gaussian scale score is
$s_{d_p}=(d_p-|Z|^2)/(2\vartheta)$, with $Z$ standard Gaussian.
It has zero mean and second moment $d_p/(2\vartheta^2)$.
The tilted density ratio lies in $[e^{-B_0},e^{B_0}]$, so
\begin{equation}
 \partial_\vartheta\psi_{p,j}=-\E_{\nu_{p,x,j}}s_{d_p},\qquad
 |\partial_\vartheta\psi_{p,j}|
 \le(e^{B_0}-1)\frac{\sqrt{d_p}}{\sqrt2\vartheta}.
 \label{eq:rg-formation-tilt-score-proof}
\end{equation}
This bound is uniform in the retained input because the Gaussian
conditional center is independent of $\vartheta$.

The full restricted objective retains its common scalar $C_T$.
Differentiating it at fixed actual $q$ gives
\begin{align}
 g_p(q,\vartheta)
 ={}&\vartheta-2+\frac12\E_qx^TH_px
       +\frac{Td_p}{2\vartheta}
       +T\E_q\partial_\vartheta\psi_{p,S}
 \notag\\
 &-\frac T2\Tr(A_0^{-1}H)
       -T\E_q\partial_\vartheta\log\pi_S(\vartheta).
 \label{eq:rg-formation-complete-score-proof}
\end{align}
The first term is the calibration derivative, and the trace is the
derivative of the proper prior's full Gaussian normalizer. The source
term is present even though it has no continuous-coordinate gradient.
For the distinguished binary prior its absolute unscaled score is at
most four. A time derivative of actual source masses is not a term in
this partial derivative.

For completeness, the complete generative Fisher is
\begin{equation}
 M(\vartheta)=T^{-1}
   +\frac12\Tr[(A_0^{-1}H)^2]
   +\operatorname{Var}_{\pi_{i_0}(\vartheta)}(t),
 \qquad t_+=4,\quad t_-=0,
 \qquad m(\vartheta)=TM(\vartheta).
 \label{eq:rg-formation-inherited-metric-proof}
\end{equation}
The prior, calibration, and source-prior scores have zero mean and
zero cross expectations in the normalized joint. Its other conditional
kernels have no explicit parameter derivative at fixed inputs. This
proves the formula without substituting a posterior or marginal
Fisher. Since $R$ is fixed, annihilates $\mathcal K$, and is supported
on the first normal layer, $A_0=\vartheta H+P_{\mathcal K}-TR$ is
positive for sufficiently small $T$. Conjugation by the positive
quadratic core bounds the trace terms and their derivatives uniformly.
Thus
\begin{equation}
 1\le m(\vartheta)\le1+C_LT,\qquad
 |m'(\vartheta)|\le C_LT.
 \label{eq:rg-formation-metric-bounds-proof}
\end{equation}
All normalizers, determinant derivatives, and categorical terms in
\cref{eq:rg-formation-complete-score-proof} follow from normalized
integrals. The physical prior correction $-TR$ has not been inserted
into the optimizing Gaussian precision. \hfill$\square$

\section{Prepared actual laws and coupled continuation}
\label{sec:rg-formation-continuation-proof}

\lemmaheading{Prepared actual-law continuation}
{lem:formation-prepared-continuation}
For the admitted common-center preparation and every finite sequence
of actual marked activations, the coupled diffusion, source proposals,
and own-law parameter equation have a unique nonexplosive solution
through $H_*=2/\epsilon$. The actual laws remain positive smooth
joint densities and satisfy the moment estimates below. The compact
parameter and speed bounds are established from those moments before
any conditional-entropy or source-history estimate is used.
\status{ESTABLISHED}

\emph{Proof.}
For a fixed positive instance and a prescribed continuous parameter
path in a compact neighborhood of $[1,2.1]$, each sheet drift has a
globally bounded Hessian and linear growth by
\cref{eq:rg-formation-parent-relative-bounds}. All proposal rows are
smooth, have bounded derivatives for this instance, and have a
positive finite-depth lower bound. Their denominator in the Gaussian
integral representation is at least $e^{-B_0}$. Interlacing the sheet
diffusions on an independent rate-one Poisson grid therefore constructs
a unique nonexplosive process with all fixed-order polynomial moments.

To close the parameter feedback, let
$u_j^{\vartheta}(s,x)=\E_{\vartheta}[f(X_t,S_t)\mid X_s=x,S_s=j]$
for a smooth terminal function whose derivatives have polynomial
growth. Condition on the proposal times and a finite label sequence.
The resulting expectation is a composition of Lipschitz sheet flows
and products of positive row weights. Flow derivatives are bounded on
compact time intervals; a differentiated weight is bounded by a
constant times that weight, since the weights have a positive lower
bound. The proposal count has every exponential moment. Polynomial
diffusion moments therefore justify differentiating this
representation and give polynomial bounds for $u$ and $\nabla u$.
Applying It\^o's formula to $u^{\widetilde\vartheta}$ along the
$\vartheta$ process, first stopped in a bounded ball and then using
these moments, proves
\begin{equation}
 |\E_{\vartheta}f(X_t,S_t)
       -\E_{\widetilde\vartheta}f(X_t,S_t)|
 \le C\int_0^t|\vartheta(s)-\widetilde\vartheta(s)|ds.
 \label{eq:rg-formation-local-law-dependence}
\end{equation}
The second-order coefficients agree, the drift difference is bounded
by $C|\vartheta-\widetilde\vartheta|(1+|x|)$, and the row difference
by $C|\vartheta-\widetilde\vartheta|$. The observable terms in
\cref{eq:rg-formation-complete-score-proof} meet these hypotheses.
Their explicit terminal parameter dependence adds a local Lipschitz
term. Hence the integral map for
$\dot\vartheta=-\epsilon m(\vartheta)^{-1}g_p(q,\vartheta)$ is a
contraction on a sufficiently short interval. This proves local
coupled existence and uniqueness. These local constants may depend on
the fixed positive instance; no uniform short-time contraction
constant is asserted.

Stop only at the first possible exit of $\vartheta$ from $[1,2.1]$.
The Gaussian comparison starts at common non-root center $Av_0$,
zero loop center, and covariance $TH_0^{-1}$. Its initial coupling
to every admitted sheet satisfies
$|H_0^{1/2}(X_0-X_0^G)|<\epsilon\sqrt T$. The root is the prescribed
independent equilibrium factor. Small compactly supported smooth
transports in standardized non-root coordinates satisfy this
displacement bound and a fixed bounded log-density-ratio bound, so the
preparation neighborhood is nonempty. In particular no order-$A$
spread of sheet centers is admitted.

On an active level let the Gaussian comparison use the same actual
deterministic parameter and Brownian increments:
$dX^G=-\vartheta(t)H_pX^Gdt+\sqrt{2T}dW$. Continue its projected
reference at an activation, while initializing the actual parent at
the full actual retained sample. For $e=X-X^G$, the noise cancels and
\cref{eq:rg-formation-parent-relative-bounds} gives
\begin{equation}
 \frac d{dt}|e|^2\le-e^TH_pe+C_LT,
 \qquad
 \sup_{t\le H_*}|e_t|
 \le |e_0|+C_L\sqrt{TH_*}.
 \label{eq:rg-formation-path-force-bound}
\end{equation}
A source proposal leaves $e$ continuous. The estimate is valid for
every adapted label history, without conditioning Brownian increments
on that history.

For the energy estimate, variation of constants uses the commuting
propagator $\exp[-H_p\int_s^t\vartheta(r)dr]$. Spectral maximization
gives
\begin{equation}
 \int_0^\infty\|H_pe^{-H_pu}\|du
 \le C_L\left[1+\log
       \frac{\lambda_{\max}(H_p)}{\lambda_{\min}(H_p)}\right].
 \label{eq:rg-formation-logarithmic-propagator}
\end{equation}
Indeed split at the inverse largest and smallest eigenvalues. The
three bounds are the largest eigenvalue, $1/(eu)$, and the smallest
eigenvalue times its decaying exponential. Their integrals give the
constant, logarithmic, and tail terms. Euclidean projection contracts
\cref{eq:rg-formation-path-force-bound}, and the Schur variational
identity contracts its energy counterpart at each activation. With
$\ell_\epsilon=1+\log(1/\epsilon)$, finitely many such stages give
\begin{equation}
 \sup_t|H_p^{1/2}e_t|\le C_L\sqrt T\ell_\epsilon.
 \label{eq:rg-formation-energy-force-bound}
\end{equation}

The Gaussian covariance barrier
$\Sigma^G\preceq C_LT H_p^{-1}$ uses only $\vartheta\ge1$.
It follows either from the stochastic integral or by substituting the
constant upper barrier into
$\dot\Sigma^G=-\vartheta H_p\Sigma^G
-\vartheta\Sigma^GH_p+2TI$.
At retention, $P H_p^{-1}P^T=H_s^{-1}$ preserves that barrier.
The deterministic mean contracts, has norm $O_L(A)$, and has zero
loop part. Combining this with
\cref{eq:rg-formation-energy-force-bound} yields
\begin{align}
 \E_{q_t^p}x^TH_px
 &\le C_L[A^2J_{p+1}+T\ell_\epsilon^2],
 \label{eq:rg-formation-actual-energy}\\
 \E_{q_t^p}|x|^2
 &\le C_L[A^2+T\ell_\epsilon^2],
 \qquad J_1\E_{q_t^p}|h|^2\le C_LT.
 \label{eq:rg-formation-actual-moments}
\end{align}
The sharper loop bound follows separately from its scalar core rate
$\vartheta J_1$ and force at most $C_L\sqrt{TJ_1}$; its forced error
is at most $C_L\sqrt{T/J_1}$. This does not assert that the nonlinear
loop itself is an independent Ornstein--Uhlenbeck process.

Set
\begin{equation}
 r_\epsilon=A^2J_1+T\ell_\epsilon^2.
 \label{eq:rg-formation-parameter-remainder}
\end{equation}
Substitution in the complete covector and metric bounds proves
$\dot\vartheta=\epsilon(2-\vartheta)
+O_L(\epsilon r_\epsilon)$ before assuming any bound on
$\dot\vartheta$. The field points inward at both endpoints of
$[1,2.1]$ for small $\epsilon$, excluding the stopped exit.
The moment bounds exclude explosion. This extends the solution
through $H_*$ and gives $|\dot\vartheta|\le C_L\epsilon$.
Positive-temperature elliptic smoothing and positive finite rates
preserve the positive smooth densities. The same construction applies
after an actual activation because its full retained law, rather than
an equilibrium replacement, is the initial law. \hfill$\square$

\lemmaheading{Comparison of independently learned parameters}
{lem:formation-parameter-comparison}
For every actual process or finite replacement lineage in
\cref{lem:formation-prepared-continuation}, the following estimates
are established:
\begin{equation}
 1\le\vartheta_p(t)\le2.1,\qquad
 |\dot\vartheta_p(t)|\le C_L\epsilon,\qquad
 \sup_{t\le H_*}|\vartheta_p(t)-(2-e^{-\epsilon t})|
 \le C_Lr_\epsilon.
 \label{eq:rg-formation-parameter-comparison-proof}
\end{equation}
Every independently evolving descendant uses its own actual law and
copies only the actual starting parameter. \status{ESTABLISHED}

\emph{Proof.}
The first two bounds were proved without entropy or averaging.
Subtract $\dot\vartheta_0=\epsilon(2-\vartheta_0)$,
$\vartheta_0(0)=1$, from the actual equation and integrate with factor
$e^{\epsilon t}$. The forcing contribution is at most $C_Lr_\epsilon$
through the full horizon. Copying the actual parameter at each of
finitely many activations adds only a finite multiple of that error.
Under the declared regime
\begin{equation}
 r_\epsilon
 =\epsilon^{2Q_1+20}
   +\epsilon^{3Q_1+26}\ell_\epsilon^2\longrightarrow0.
 \label{eq:rg-formation-small-parameter-remainder}
\end{equation}
In particular each coefficient exceeds $7/4$ on the last ten units of
the horizon, for sufficiently small $\epsilon$, since
$2-e^{-2}>7/4$. The comparison trajectory is a consequence of the
own-law equation, not an imposed coefficient schedule.
\hfill$\square$

\section{Entropy estimates for actual joint laws}
\label{sec:rg-formation-entropy-proof}

\lemmaheading{Continuous entropy and score tools}
{lem:formation-entropy-tools}
Let $\mu(dx)\propto e^{-W(x)/T}dx$ on a finite-dimensional Euclidean
space, where $W$ is smooth, has a globally bounded Hessian, and
$\nabla^2W\succeq\lambda I$ with $\lambda>0$. For a probability density
$f$ relative to $\mu$, the following inequalities hold whenever their
right sides are finite:
\begin{align}
 \operatorname{Ent}_\mu(f)
 &\le \frac{T}{2\lambda}
           \int\frac{|\nabla f|^2}{f}d\mu,
 \label{eq:rg-formation-continuous-lsi}\\
 \log\E_\mu e^{u(F-\E_\mu F)}
 &\le\frac{T K^2u^2}{2\lambda},
 \qquad \operatorname{Lip}(F)\le K.
 \label{eq:rg-formation-herbst-proof}
\end{align}
Here $\operatorname{Ent}_\mu(f)=\int f\log f d\mu$ when $\int f d\mu=1$;
the unnormalized convention used below is
$\operatorname{Ent}_\mu(f)=\int f\log(f/\int f d\mu)d\mu$.
These statements apply to every conditional sheet potential in this
construction. \status{ESTABLISHED}

\emph{Proof.}
For the semigroup $P_t$ generated by $T\Delta-\nabla W\cdot\nabla$,
the derivative of the synchronous stochastic flow has operator norm
at most $e^{-\lambda t}$. Weighted Cauchy--Schwarz therefore gives
$|\nabla P_tf|^2/P_tf\le e^{-2\lambda t}
P_t(|\nabla f|^2/f)$. Invariance of $\mu$ implies
$I_\mu(P_tf)\le e^{-2\lambda t}I_\mu(f)$.
Integration by parts gives
$d\operatorname{Ent}_\mu(P_tf)/dt=-TI_\mu(P_tf)$.
For bounded positive Lipschitz $f$, contraction against a stationary
initial sample yields $P_tf(x)\to\int f d\mu$. Bounded convergence
then makes its entropy tend to zero. Integrating the Fisher bound
proves \cref{eq:rg-formation-continuous-lsi} with its stated constant.
Cutoff, positive truncation, and short semigroup smoothing extend it
by lower semicontinuity. Gaussian tails from strong convexity justify
the integrations. This is the curvature entropy argument in
\citet[pp.~187, 194--200]{BakryEmery1985}, with diffusion coefficient
$T$ made explicit.

Apply \cref{eq:rg-formation-continuous-lsi} to $e^{uF}$.
For $K_F(u)=\log\E_\mu e^{uF}$ it gives
$uK_F'(u)-K_F(u)\le TK^2u^2/(2\lambda)$.
Integration of the derivative of $K_F(u)/u$ from zero proves
\cref{eq:rg-formation-herbst-proof}, also for negative $u$.
The entropy variational inequality
\begin{equation}
 u\E_\nu F\le D(\nu\Vert\mu)+\log\E_\mu e^{uF}
 \label{eq:rg-formation-entropy-duality}
\end{equation}
follows from nonnegativity of relative entropy to the normalized
exponential tilt of $\mu$. Consequently
$|\E_\nu F-\E_\mu F|\le K\sqrt{2TD(\nu\Vert\mu)/\lambda}$.
If instead a centered observable $\xi$ obeys
\begin{equation}
 \log\E_\mu e^{u\xi}\le
       \frac{v u^2}{2(1-b|u|)},\qquad |u|<b^{-1},
 \label{eq:rg-formation-subgamma-proof}
\end{equation}
optimization in \cref{eq:rg-formation-entropy-duality}, using
$|u|=\sqrt{2D/v}/(1+b\sqrt{2D/v})$, gives
\begin{equation}
 |\E_\nu\xi|
 \le\sqrt{2vD(\nu\Vert\mu)}+bD(\nu\Vert\mu).
 \label{eq:rg-formation-score-duality}
\end{equation}
The $D=0$ case follows by continuity. Finally
$\|\nu-\mu\|_{\mathrm{TV}}\le\sqrt{D(\nu\Vert\mu)/2}$:
push both laws to a set realizing their total variation, and use
the binary relative entropy's second derivative
$[t(1-t)]^{-1}\ge4$ and its zero value and derivative at equality.
\hfill$\square$

All normalized hidden Gaussian conditionals have a center independent
of $\vartheta$. Their logarithmic parameter scores, after including
the bounded marginal tilt, are centered versions of
$-|Z|^2/(2\vartheta)-\partial_\vartheta\psi_{p,j}$.
The full joint score adds
$\partial_\vartheta\log\pi_j(\vartheta)$, whose absolute value is at
most four. These scores satisfy
\cref{eq:rg-formation-subgamma-proof} with constants $v_L,b_L$
independent of the retained input and of $T,J_s$. To check this
uniformity directly, a bounded tilt of a standard Gaussian satisfies
\begin{equation}
 \E|Z|^{2m}
 \le e^{B_0}2^{k_0/2}4^m m!,\qquad m\ge0.
 \label{eq:rg-formation-gaussian-score-moments}
\end{equation}
Indeed $r^m\le m!4^m e^{r/4}$ and the Gaussian exponential moment is
$2^{k_0/2}$. The other score terms are uniformly bounded by
\cref{eq:rg-formation-tilt-score-proof}. Centering increases a
suitable geometric moment constant by at most a factor two. Expanding
the exponential series, removing the zero linear term, and using
$\log(1+x)\le x$ proves \cref{eq:rg-formation-subgamma-proof}.
Conditional expectations of such scores retain the same bound by
Jensen's inequality. \status{ESTABLISHED}

\lemmaheading{Entropy at actual activations}
{lem:formation-activation-entropy}
For the actual active law at level $p$, let
$E_p(t)=T D(q_t^p\Vert\pi_{\vartheta_p(t)}^p)$, with the independent
root factor canceled. At an activation time $\tau_p$,
\begin{equation}
 E_p(\tau_p)\le C_L[A^2J_{p+1}+T\ell_\epsilon^2].
 \label{eq:rg-formation-activation-budget}
\end{equation}
On each ensuing active interval,
\begin{equation}
 T^2\int_{\tau_p}^{t}I_p(u)du
 \le C_L[E_p(\tau_p)+T],
 \label{eq:rg-formation-full-fisher-budget}
\end{equation}
where $I_p$ is the full continuous relative Fisher information of
the joint actual law. Neither estimate assumes a positive lower
bound on its time-dependent source masses. \status{ESTABLISHED}

\emph{Proof.}
We first control actual differential/counting entropy
$h(q)=-\sum_j\int q_j\log q_j$ across changes of dimension.
Write $q_j=\lambda_jf_j$, with terms of zero mass omitted. On a
$d$-dimensional sheet, diffusion gives
\begin{equation}
 h'(q)_{\mathrm{diff}}
 =T\sum_j\lambda_jI(f_j)-\E_q\Delta V_{p,S}
 \ge T\sum_j\lambda_jI(f_j)-d\Lambda_p,
 \qquad \Lambda_p\le C_LJ_1.
 \label{eq:rg-formation-actual-entropy-derivative}
\end{equation}
The entropy--Fisher bound
$I(f)\ge2\pi e d e^{-2h(f)/d}$ follows from the Gaussian case of
\cref{eq:rg-formation-continuous-lsi}: expand both sides relative to
$N(0,\sigma^2I)$, integrate the cross term
$\int x\cdot\nabla f=-d$, cancel the second moments, and minimize
at $\sigma^2=d/I(f)$. Jensen's inequality and
$h(q)=\sum_j\lambda_jh(f_j)+H(\lambda)$, with $H(\lambda)\ge0$,
then give
$\sum_j\lambda_j I(f_j)\ge2\pi e d e^{-2h(q)/d}$.

At fixed $x$, a proposal moves the conditional label row $\alpha$
in the direction $\alpha R-\alpha$. Concavity of categorical
entropy yields its directional derivative at least
$H(\alpha R)-H(\alpha)\ge-\log|\mathcal J|$. The continuous
marginal is unchanged by the proposal. Hence
\begin{equation}
 h'(q)\ge 2\pi e d T e^{-2h(q)/d}
                -d\Lambda_p-\log|\mathcal J|.
 \label{eq:rg-formation-counting-entropy-floor}
\end{equation}
Scalar comparison bounds $h(q)$ below by the lesser of its initial
value and
$\frac d2\log[2\pi e dT/(d\Lambda_p+\log|\mathcal J|)]$.
At a projection dropping $k$ coordinates $z$, the chain rule and
the maximum-entropy Gaussian bound imply
\begin{equation}
 h(Y,S)\ge h(X,S)
       -\frac k2\log\frac{2\pi e \E|z-z_*|^2}{k}.
 \label{eq:rg-formation-entropy-projection}
\end{equation}
Here $z_*$ can be the projected comparison mean. The bound follows
conditionally on $(Y,S)$ and then from Jensen; it also follows from
nonnegativity of relative entropy to the common isotropic Gaussian.
The moment proof of
\cref{lem:formation-prepared-continuation} bounds this centered
covariance by $C_LT\ell_\epsilon^2H_p^{-1}$. Initial entropies and
all logarithms in these estimates are $O_L(\ell_\epsilon^2)$:
the covariance eigenvalues, $A$, and $T$ are fixed powers of
$\epsilon$, and the admitted initial density ratios are bounded.
There are finitely many projections. Expanding
$TD(q\Vert G_\vartheta^p\pi(\vartheta))$ and using
\cref{eq:rg-formation-actual-energy} proves
\cref{eq:rg-formation-activation-budget}. The bounded optimizer
tilt adds at most $TB_0$. The source prior has a positive minimum
on $[1,2.1]$, so its log contribution is bounded independently of
the actual $\lambda_j$.

For the time-integrated estimate put $V=x^TH_px/2$.
The normalized reference satisfies
\begin{equation}
 \E_{\pi_\vartheta^p}\exp\frac{\vartheta V}{2T}
 \le R_L,\qquad R_L=e^{B_0}2^{k_0/2}.
 \label{eq:rg-formation-energy-exponential-moment}
\end{equation}
Entropy duality gives
$\vartheta\E_qV\le2E_p+2T\log R_L$; Jensen gives
$\vartheta\E_\pi V\le2T\log R_L$. In the normalized full score,
the remaining centered terms have a uniform bound $2M_L$.
Therefore
\begin{equation}
 |T\E_q\partial_\vartheta\log\pi_\vartheta^p|
 \le\frac{2E_p+4T\log R_L}{\vartheta}+2TM_L.
 \label{eq:rg-formation-entropy-parameter-budget}
\end{equation}
Diffusion dissipates $T^2I_p$. Detailed balance
\cref{eq:rg-formation-parent-balance-proof} makes the source
contribution nonpositive. Thus, for a fixed $C_{0,L}$,
\begin{equation}
 E_p'\le -T^2I_p+
    \frac{2|\dot\vartheta|}{\vartheta}(E_p+C_{0,L}T).
 \label{eq:rg-formation-entropy-energy-inequality}
\end{equation}
The integrating factor is
$\exp[-2\int_{\tau_p}^t|\dot\vartheta|/\vartheta]$.
Its exponent is bounded over $H_*$ by the already proved speed
bound. Integration proves \cref{eq:rg-formation-full-fisher-budget}.
At an exact activation both actual and reference laws are pushed by
the same channel at the same current parameter. The relative-entropy
chain rule makes $E_p$ decrease, so the estimates continue through
the declared finite lineage. All entropy differentiations follow
first for positive truncated densities on bounded domains; the
moment and entropy bounds above justify removing these cutoffs.
\hfill$\square$

\lemmaheading{Conditional entropy and complete score error}
{lem:formation-conditional-entropy}
For every retention $p=s-1\to s$, the actual conditional relative
entropy with respect to the optimizing hidden conditional satisfies
\begin{equation}
 \sup_{t\in I_s\ \mathrm{and\ subsequent\ active\ times}}
 \Delta_s(t)
 \le
 \begin{cases}
 C_L\epsilon^{14},&s<L,\\
 C_L\epsilon^2,&s=L.
 \end{cases}
 \label{eq:rg-formation-conditional-final-bounds}
\end{equation}
The early bound also applies until the next activation to the
corresponding un-replaced comparison process. At every such time
the complete fine/parent covector difference is at most
$C_LT(\sqrt{\Delta_s}+\Delta_s)$.
The terminal assertion is proved below using joint relaxation.
\status{ESTABLISHED}

\emph{Proof of the early bound and the covector estimate.}
For an early retention use the auxiliary continuous split
\begin{equation}
 Z=(\mathcal N_p\cap\mathcal N_s^\perp)\oplus\mathcal H,
 \qquad Y=\mathcal N_s,
 \qquad S\ \text{retained}.
 \label{eq:rg-formation-auxiliary-split}
\end{equation}
The true parent retains the loops. The entropy chain rule therefore
bounds its hidden conditional defect $\Delta_s$ by the auxiliary
defect $\widehat\Delta_s$ obtained by also integrating $h$.
The independent root cancels. Conditional strong convexity in $Z$
is at least $\lambda_s=c_{\mathrm{early},L}J_s$ after reducing the
fixed positive constant if necessary. The mixed Hessian satisfies
\begin{equation}
 \|\nabla_Z\nabla_YV_{p,j}\|\le\Gamma_s=C_LJ_{s+1},
 \qquad \Gamma_L=0.
 \label{eq:rg-formation-mixed-hessian-scale}
\end{equation}
Here is the scale check, which a bound by $J_1$ would not replace.
The only nonlinear residuals lie in the first normal layer and the
loop space. At $p=0$ the original nonlinear potential has no future
normal dependence, so its mixed future derivatives vanish directly.
For $p\ge1$, use coordinates $r_E=B_Ew_E$ on eliminated normal rows.
The Gaussian conditional quadratic has precision
$A_E=D_E+C^TD_fC$. Its perturbation
$\mathcal R_E=C^TD_fC$ has norm at most $C_LJ_{p+1}$,
while $A_E^{-1}$ has norm at most $C_L/J_p$.
The identity
\begin{equation}
 A_E^{-1}=D_E^{-1}
           -D_E^{-1}\mathcal R_EA_E^{-1}
 \label{eq:rg-formation-first-row-resolvent}
\end{equation}
bounds the first block row by $C_L/J_1$.
A future-layer $r$ force column has size $C_LJ_r$, so the first
conditional mean responds by $C_LJ_r/J_1$. Its conditional
covariance is comparable to $T/J_1$.

Differentiate the normalized Gaussian integral defining
$\psi_{p,j}$. A first-layer or loop derivative of $\phi_j$ has
size $C_L\sqrt{J_1/T}$, and a second derivative has size
$C_LJ_1/T$. Each derivative in a future direction brings the
conditional-mean response just obtained. The covariance term in the
second derivative is bounded by the product of the same first
derivative bounds. Consequently
\begin{equation}
 T|\partial_h\partial_{w_r}\psi_{p,j}|\le C_LJ_r,\qquad
 T|\partial_{w_i}\partial_{w_r}\psi_{p,j}|
       \le C_L\frac{J_iJ_r}{J_1}.
 \label{eq:rg-formation-future-derivative-scales}
\end{equation}
The Gaussian Schur cross block has size at most $C_LJ_{s+1}$,
by the lower triangular residual form and its variational bound.
Summing finitely many blocks proves
\cref{eq:rg-formation-mixed-hessian-scale}. This calculation includes
the generated conditional covariance.

Write the joint density as $q=r(y,j)a(z\mid y,j)$ and its
optimizing reference as $\pi=\pi_Y(y,j)k(z\mid y,j)$.
Let
$w_Y=\nabla_Y\log(q/\pi)$,
$v_Y=\nabla_Y\log(r/\pi_Y)$, and
$e_Y=\E_a\nabla_Y\log k$. Normalization gives
$\E_k\nabla_Y\log k=0$ and $\E_aw_Y=v_Y-e_Y$.
The actual marginal flux is $Tr(v_Y-e_Y)$.
Subtracting its entropy derivative from the full derivative yields
the exact continuous identity
\begin{align}
 \widehat\Delta_s'|_{\mathrm{cont}}
 ={}&-TI_Z
      -T\E_r\operatorname{Var}_a(w_Y)
      +T\E_r[e_Y\cdot\E_aw_Y]
 \notag\\
 &-\dot\vartheta \E_q\partial_\vartheta\log k.
 \label{eq:rg-formation-conditional-balance}
\end{align}
The conditional logarithmic Sobolev inequality gives
$TI_Z\ge2\lambda_s\widehat\Delta_s$.
Apply \cref{eq:rg-formation-herbst-proof} to each unit-vector
component of $\nabla_Y\log k$, whose $Z$ Lipschitz constant is at
most $\Gamma_s/T$. This gives
$|e_Y|\le\Gamma_s\sqrt{2D(a\Vert k)/(T\lambda_s)}$.
Cauchy--Schwarz and Young's inequality bound the mixed term in
\cref{eq:rg-formation-conditional-balance} by
$\lambda_s\widehat\Delta_s/2+
T\Gamma_s^2I_Y/\lambda_s^2$; $I_Y$ is the actual full relative
Fisher component. The normalized conditional parameter score has
the uniform exponential-moment bound above. Its contribution is at
most
$\lambda_s\widehat\Delta_s/2+
v_L|\dot\vartheta|^2/\lambda_s+
b_L|\dot\vartheta|\widehat\Delta_s$.

A finite-rate source update is not a conditional-entropy contraction,
so its contribution must be retained. At fixed $Y=y$, compare $k_j$
with the common label-independent hidden Gaussian $\gamma$.
Each $\log(\gamma/k_j)$ lies in $[-B_0,B_0]$.
The conditional entropy relative to $\gamma$ splits as
$D(q_Z\Vert\gamma)+I(Z;S)$.
A proposal leaves $q_Z$ unchanged. Mutual information is convex as
a function of the conditional label channel at fixed $q_Z$, by joint
convexity of relative entropy. Its directional derivative towards
the new channel is at most $\log|\mathcal J|$.
The difference from the actual target is a linear expectation of
$\log(\gamma/k_j)$ and changes at rate at most $2B_0$.
Thus
\begin{equation}
 \widehat\Delta_s'|_{\mathrm{jump}}\le
 C_J:=\log|\mathcal J|+2B_0.
 \label{eq:rg-formation-conditional-jump-remainder}
\end{equation}
This argument remains valid when a source mass vanishes, by convex
approximation, and divides by no such mass.

For sufficiently small $\epsilon$,
$b_L|\dot\vartheta|\le\lambda_s/2$. Dropping the negative variance
term gives
\begin{equation}
 \widehat\Delta_s'
 \le-\frac{\lambda_s}{2}\widehat\Delta_s
     +\frac{T\Gamma_s^2}{\lambda_s^2}I_Y
     +\frac{v_L|\dot\vartheta|^2}{\lambda_s}
     +C_J.
 \label{eq:rg-formation-conditional-differential-bound}
\end{equation}
The initial conditional entropy is at most $E_p(\tau_p)/T$.
Integrate the inequality and apply
\cref{eq:rg-formation-full-fisher-budget,eq:rg-formation-activation-budget}.
Since $A^2/T=\epsilon^{-6}$, the result is
\begin{align}
 \widehat\Delta_s(t)\le C_L\bigg[
 &e^{-c_{\mathrm{early},L}J_s(t-\tau_p)}
                 (\epsilon^{-6}J_s+\ell_\epsilon^2)
 \notag\\
 &+\epsilon^{-6}\frac{J_{s+1}^2}{J_s}
 +\ell_\epsilon^2\frac{J_{s+1}^2}{J_s^2}
 +\frac{\epsilon^2}{J_s^2}+\frac{C_J}{J_s}\bigg],
 \label{eq:rg-formation-early-entropy-bound}
\end{align}
where the constant named $c_{\mathrm{early},L}$ is chosen to include
the preceding factor $1/2$. For $s<L$,
$J_{s+1}^2/J_s=\epsilon^{20}$ and $J_s\ge\epsilon^{-20}$.
The declared entrance time and
$c_{\mathrm{early},L}k_L/2>Q_1+20$ make the first term no larger
than $C_L\epsilon^{14}$. The other terms meet the same bound.
This proves the early assertion. The last term would be $C_J$ at
$s=L$; it is not discarded there.

For the true fixed retention the exact objective chain rule is
\begin{equation}
 F_T^p(q,\vartheta)
 -F_T^s((C_{s\leftarrow p})_\#q,\vartheta)
 =T\Delta_s(q,\vartheta).
 \label{eq:rg-formation-objective-chain-proof}
\end{equation}
Differentiating at fixed actual $q$ and fixed channel gives
$g_p-g_s=-T\E_q\partial_\vartheta\log k_{\mathrm{true}}$.
The normalized conditional score has the same uniform
exponential-moment bound. Apply
\cref{eq:rg-formation-score-duality} conditionally and then Jensen
to obtain
\begin{equation}
 |g_p-g_s|
 \le C_LT(\sqrt{\Delta_s}+\Delta_s).
 \label{eq:rg-formation-complete-score-error}
\end{equation}
The common scalar and source-prior derivative cancel through this
exact chain rule; they were not removed from either objective.
\hfill$\square$

\lemmaheading{Joint relaxation at the fixed weakest strength}
{lem:formation-terminal-joint-entropy}
There is a constant $c_{\mathrm{joint},L}>0$, independent of
$\epsilon$, such that every actual lineage satisfies
\begin{equation}
 D(q_t^p\Vert\pi_{\vartheta_p(t)}^p)
 \le C_L e^{-c_{\mathrm{joint},L}t}
           [\epsilon^{-(Q_1+6)}+\ell_\epsilon^2]
       +C_L\epsilon^2.
 \label{eq:rg-formation-terminal-joint-bound}
\end{equation}
Consequently the terminal part of
\cref{eq:rg-formation-conditional-final-bounds} holds once
$t\ge t_L=k_L\log(1/\epsilon)$. \status{ESTABLISHED}

\emph{Proof.}
Fix a level and parameter, write
$p=\pi_\vartheta^p$, $\rho=G_\vartheta^p\pi(\vartheta)$, and
$u_j=q_j/p_j$. The bounded tilt gives
$e^{-B_0}\le p/\rho\le e^{B_0}$, and
\cref{eq:rg-formation-inherited-row-proof} gives
\begin{equation}
 p_j(x)R_{jk}^p(x)
 \ge e^{-2B_0}G_\vartheta^p(x)
                    \pi_j(\vartheta)\pi_k(\vartheta).
 \label{eq:rg-formation-conductance-comparison}
\end{equation}
For a probability measure $\mu$,
$\operatorname{Ent}_\mu(u)=\inf_{r>0}\int[u\log(u/r)-u+r]d\mu$.
The integrand is nonnegative, so
$\operatorname{Ent}_p(u)\le e^{B_0}\operatorname{Ent}_\rho(u)$.
The product chain rule and convexity imply
\begin{equation}
 \operatorname{Ent}_\rho(u)
 \le\sum_j\pi_j\operatorname{Ent}_G(u_j)+\E_G\operatorname{Ent}_{\pi}(u(x)).
 \label{eq:rg-formation-product-entropy-bound}
\end{equation}
The Gaussian entropy inequality bounds the first term by
$T/(2\vartheta h_L)$ times its continuous Fisher information.
For any positive finite vector $v$,
\begin{align}
 \frac12\sum_{j,k}\pi_j\pi_k
       (v_j-v_k)(\log v_j-\log v_k)
 &=\sum_j\pi_jv_j\log v_j
       -\bar v\sum_j\pi_j\log v_j
 \notag\\
 &\ge \operatorname{Ent}_\pi(v),\qquad \bar v=\sum_j\pi_jv_j.
 \label{eq:rg-formation-categorical-entropy-bound}
\end{align}
The inequality is Jensen's inequality for $\log$.
Compare the Fisher measures and the conductances using their positive
pointwise bounds. For
\begin{equation}
 \mathcal D_{\mathrm{jump}}
 =\frac12\int\sum_{j,k}p_jR_{jk}^p
           (u_j-u_k)(\log u_j-\log u_k)dx,
 \label{eq:rg-formation-jump-dissipation}
\end{equation}
one obtains
\begin{equation}
 D(q\Vert p)\le
 e^{2B_0}\frac{T}{2\vartheta h_L}I_p
       +e^{3B_0}\mathcal D_{\mathrm{jump}}.
 \label{eq:rg-formation-joint-coercivity}
\end{equation}
No lower bound on the actual label masses, nor a vanishing source
time scale, occurs in this inequality. Let
$c_{0,L}=\min\{2h_Le^{-2B_0},e^{-3B_0}\}>0$.
Then $TI_p+\mathcal D_{\mathrm{jump}}\ge c_{0,L}D$.
The time-dependent normalized full score satisfies
\cref{eq:rg-formation-subgamma-proof}. Hence
\begin{equation}
 D'\le-c_{0,L}D+
       C_L\epsilon(\sqrt D+D)
 \le-\frac{c_{0,L}}{2}D+C_L\epsilon^2
 \label{eq:rg-formation-joint-moving-reference}
\end{equation}
for sufficiently small $\epsilon$.
Take $c_{\mathrm{joint},L}=c_{0,L}/2$. Actual activations decrease
$D$ at the unchanged parameter; each descendant then satisfies the
same inequality with its own parameter speed.
The initial entropy estimate is
$D_0\le C_L[\epsilon^{-(Q_1+6)}+\ell_\epsilon^2]$.
Integrating through finitely many such decreases proves
\cref{eq:rg-formation-terminal-joint-bound}.
The condition $c_{\mathrm{joint},L}k_L>Q_1+8$ makes its decaying
term $O_L(\epsilon^2)$ at $t_L$.
Both the auxiliary and true hidden conditional entropies are
bounded by full joint relative entropy. This proves the terminal
assertion without treating the nonzero conditional jump remainder
as small. \hfill$\square$

\section{Exact fast marginals and complete source histories}
\label{sec:rg-formation-history-proof}

\lemmaheading{Common exact fast marginal and environmental error}
{lem:formation-fast-marginal}
There is a normalized fast law, conditional on each assignment,
whose averaged proposal matrix is common to every exact optimizing
parent at a fixed $\vartheta$. Replacing the actual environment by
this auxiliary fast process incurs the explicit errors below; the
common matrix is an exact conditional integral, not an independently
specified source dynamics. \status{ESTABLISHED}

\emph{Proof.}
Use the fixed orthonormal normal flag from
\cref{eq:rg-formation-weighted-proof}. In these coordinates
\begin{equation}
 w^TH_{\mathcal N}w
 =J_1|\mathsf B_{11}w_1|^2
       +\sum_{s=2}^LJ_s
             \left|\sum_{r\le s}\mathsf B_{sr}w_r\right|^2.
 \label{eq:rg-formation-triangular-fast-quadratic}
\end{equation}
The nonlinear factor $\phi_j$ depends only on $(w_1,h)$.
For fixed $w_1,h$, change all later $w_s$ to residuals
$r_s=\sum_{r\le s}\mathsf B_{sr}w_r$.
The domain remains the whole Euclidean product and the Jacobian is
the constant $\prod_{s=2}^L|\det\mathsf B_{ss}|$.
Integrating those Gaussian residuals contributes a scalar independent
of $w_1,h,j$. Set $A_*=J_1\mathsf B_{11}^T\mathsf B_{11}$.
The exact fast optimizing law conditional on $S=j$ is therefore
\begin{equation}
 \mu_{\vartheta,j}^*(dw_1,dh)
 =\frac{1}{Z_{\vartheta,j}^*}
   \exp\left\{-\frac{\vartheta}{2T}
         (w_1^TA_*w_1+J_1|h|^2)-\phi_j(w_1,h)\right\}dw_1dh.
 \label{eq:rg-formation-fast-law}
\end{equation}
Write $\mu_{\vartheta,j}^{*,h}$ for its loop marginal and
$\kappa_{\vartheta,j}^*(dw_1\mid h)$ for its first-normal
conditional. Define
\begin{align}
 R_{jk}^{*,h}(h,\vartheta)
 &=\int b_k(w_1,h,\vartheta)
                   \kappa_{\vartheta,j}^*(dw_1\mid h),
 \label{eq:rg-formation-fast-loop-row}\\
 \mathcal B_{jk}(\vartheta)
 &=\int b_k(w_1,h,\vartheta)\mu_{\vartheta,j}^*(dw_1,dh)
   =\int R_{jk}^{*,h}(h,\vartheta)
                    \mu_{\vartheta,j}^{*,h}(dh).
 \label{eq:rg-formation-common-averaged-row}
\end{align}
The prior factor $\pi_j(\vartheta)$ cancels on conditioning on $j$.
The Gaussian change of variables and the tower property show that
averaging the inherited row $R_{jk}^p$ against the exact optimizing
law conditional on $j$ gives this same $\mathcal B_{jk}$ for every
$p$. In particular $\mathcal B$ need not have invariant assignment
law equal to the original prior.

For an intermediate $p\ge1$, keep the future normal coordinates
$y=(w_{p+1},\ldots,w_L)$ fixed and integrate $w_2,\ldots,w_p$.
The unperturbed conditional first-normal law is
\begin{equation}
 g_{p,\vartheta,y}
 =N\left(m_p(y),\frac T\vartheta\mathsf C_p^{-1}\right),\qquad
 A_*\preceq\mathsf C_p\preceq A_*+C_LJ_2I,\qquad
 |m_p(y)|\le C_L\frac{J_2}{J_1}|y|.
 \label{eq:rg-formation-first-normal-conditional}
\end{equation}
Here $\mathsf C_p$ is a precision matrix, distinct from the Gaussian
integration scalar $C_p^G$.
For the lower precision bound, minimize the sum of nonnegative
later residual squares with $y=0$. For the upper bound, set all
eliminated later coordinates to zero and use
$J_s\le J_2$ for $s\ge2$. At $p=L$, integration of every later
residual gives $\mathsf C_L=A_*$ and $m_L=0$.
The mean bound is the first-block-row estimate
\cref{eq:rg-formation-first-row-resolvent}; each future coordinate
of level $s$ contributes at most $C_LJ_s/J_1$.

Put
\begin{equation}
 \rho_\epsilon=\frac{J_2}{J_1},\qquad
 \gamma_\epsilon=\frac{J_2}{\sqrt{J_1T}},\qquad
 d_\epsilon(y)=\rho_\epsilon+\gamma_\epsilon|y|.
 \label{eq:rg-formation-environment-scales}
\end{equation}
The Gaussian divergence from
$g_{p,\vartheta,y}$ to
$g_\vartheta^*=N(0,(T/\vartheta)A_*^{-1})$ is
\begin{equation}
 \frac12\left[
  \frac{\vartheta}{T}m_p^TA_*m_p
  +\Tr(A_*\mathsf C_p^{-1})-\dim w_1
            -\log\det(A_*\mathsf C_p^{-1})\right]
 \le C_L[\rho_\epsilon^2+\gamma_\epsilon^2|y|^2].
 \label{eq:rg-formation-first-normal-gaussian-kl}
\end{equation}
Indeed the eigenvalues of the whitened covariance lie between
$(1+C_L\rho_\epsilon)^{-1}$ and $1$, and
$t-1-\log t=O((t-1)^2)$ uniformly there for small $\epsilon$.
Pinsker's inequality gives a total-variation error at most
$C_Ld_\epsilon(y)$.

Tilting two probability measures by the same weight
$w\in[e^{-B_0},1]$ increases their total variation by at most
$2e^{B_0}$. To see this, subtract $w\mu/\mu(w)$ and
$w\nu/\nu(w)$, split numerator and normalization differences, and
use $\mu(w),\nu(w)\ge e^{-B_0}$.
Apply this to $w=e^{-\phi_j(w_1,h)}$, uniformly for every $h$.
The exact parent and fast-reference rows and loop forces therefore
satisfy
\begin{align}
 \|R_{j\cdot}^p(y,h,\vartheta)
          -R_{j\cdot}^{*,h}(h,\vartheta)\|_{\mathrm{TV}}
 &\le C_Ld_\epsilon(y),
 \label{eq:rg-formation-parent-row-environment}\\
 |\nabla_hV_{p,j}(y,h,\vartheta)
           -\nabla_hV_j^{*,h}(h,\vartheta)|
 &\le C_L\sqrt{TJ_1} d_\epsilon(y).
 \label{eq:rg-formation-parent-force-environment}
\end{align}
The second formula uses the common Gaussian force
$\vartheta J_1h$ and
$T\E_\kappa\nabla_h\phi_j$.
Furthermore
\begin{equation}
 \operatorname{Lip}_{(w_1,h)}b_{j\cdot}
       +\operatorname{Lip}_h R_{j\cdot}^{*,h}
 \le C_L\sqrt{J_1/T}.
 \label{eq:rg-formation-row-fast-lipschitz}
\end{equation}
The derivative of $R^{*,h}$ includes the covariance of $b_k$ with
$\nabla_h\phi_j$; the conditional Gaussian $g_\vartheta^*$ is
independent of $h$. Its bounded source derivatives and this
covariance prove the displayed estimate. The exact marginal
covariance has therefore been included in both the force and row
calculations. \hfill$\square$

\lemmaheading{Fast mixing with a learned coefficient}
{lem:formation-fast-mixing}
Let the auxiliary process have, on a proposal-free interval, fixed
assignment $j$ and either the full fast potential of
\cref{eq:rg-formation-fast-law} or its exact loop marginal.
Its actual deterministic coefficient satisfies
$1\le\vartheta\le2.1$ and $|\dot\vartheta|\le C_L\epsilon$.
For a point start $z$ at time $s$ and
$t-s\ge(c_{\mathrm{fast},L}J_1)^{-1}$,
\begin{equation}
 \|P^*_{s,t}(z,\cdot)-\mu^*_{\vartheta(t),j}\|_{\mathrm{TV}}
 \le C_L(1+\sqrt{J_1/T}|z|)
              e^{-c_{\mathrm{fast},L}J_1(t-s)}
       +C_L\epsilon/J_1.
 \label{eq:rg-formation-fast-mixing-bound}
\end{equation}
The notation $\mu^*$ in the loop case means its loop marginal.
Also $\mathcal B_{jk}(\vartheta)$ is uniformly Lipschitz in
$\vartheta$. \status{ESTABLISHED}

\emph{Proof.}
The relative Hessian bounds and the normalized marginal covariance
calculation give
$c_LJ_1I\preceq\nabla^2V^*\preceq C_LJ_1I$ for both potentials.
Both are even, minimize at zero, and their equilibrium second
moments are at most $C_LT/J_1$.
For two point starts, reflection coupling before meeting gives their
distance a one-dimensional noise of size $2\sqrt{2T}$ and a drift
at most $-\lambda D$, where $\lambda=c_LJ_1$.
Multiplying the comparison equation by $e^{\lambda(t-s)}$ gives
a time-changed Brownian motion of variance
$4T(e^{2\lambda(t-s)}-1)/\lambda$.
The reflection principle therefore bounds the probability of
nonmeeting by
\begin{equation}
 C|z-z'|\sqrt{\lambda/T}
                 (e^{2\lambda(t-s)}-1)^{-1/2}.
 \label{eq:rg-formation-reflection-bound}
\end{equation}
This construction uses an orthogonal predictable reflection of one
Brownian motion, so both marginals have the required Brownian law.
The drifts are globally Lipschitz for each fixed instance; construct
away from the diagonal, then join the copies at meeting. The
one-dimensional comparison and the displayed hitting bound follow
by stopping before the diagonal and taking the stopping radius to
zero. Averaging $z'$ against the equilibrium at time $s$ gives the
first term in \cref{eq:rg-formation-fast-mixing-bound}.

It remains to compare that evolving equilibrium start with the
instantaneous equilibrium at time $t$. After standardizing
$(w_1,h)$ by $\sqrt{J_1/T}$, its normalized parameter score is a
centered quadratic with uniform exponential moments.
The loop-marginal score is its conditional expectation and has the
same bound by Jensen. The moving-reference entropy calculation,
using the continuous entropy inequality with curvature $c_LJ_1$,
gives
\begin{equation}
 D'\le-c_LJ_1D+C_L\epsilon^2/J_1,\qquad D(s)=0.
 \label{eq:rg-formation-fast-equilibrium-lag}
\end{equation}
The linear-in-$D$ parameter term is absorbed for small $\epsilon$.
It follows that $D\le C_L\epsilon^2/J_1^2$, and Pinsker gives the
second term in \cref{eq:rg-formation-fast-mixing-bound}.
Finally, differentiating
$\mathcal B_{jk}=\E_{\mu_{\vartheta,j}^*}b_k$ gives the explicit
row derivative and a covariance with this centered scale score.
The sum of absolute explicit row derivatives is bounded by eight,
and the standardized score has bounded first moment. Thus
$\sum_k|\partial_\vartheta\mathcal B_{jk}|\le C_L$.
\hfill$\square$

\lemmaheading{Coupling of complete source histories}
{lem:formation-history-coupling}
The actual source-assignment histories of a process and of an
activated exact parent can be coupled, together with their common
actual initial marked sample, so that the probability of any
disagreement through $H_*$ is at most $C_L\epsilon^6$.
The result applies through the finite activation lineage. It concerns
the full finite-state path, including arbitrary actual label
correlations, and uses the declared rate-one proposal clock.
\status{ESTABLISHED}

\emph{Proof.}
We first bound an actual proposal row on one proposal-free interval.
For a microscopic process let $\xi=(w_1,h)$. Couple its auxiliary
fast diffusion to the same initial fast point, Brownian increments,
assignment, and actual parameter path. The omitted Gaussian
environmental force is
$\vartheta[(H_{11}-A_*)w_1+H_{1,>1}w_{>1}]$ and has norm at most
$C_LJ_2|w|$. Strong contraction gives the pathwise bound
\begin{equation}
 |\xi_t-\xi_t^*|
 \le C_LJ_2\int_\sigma^t
             e^{-c_LJ_1(t-v)}|w_v|dv.
 \label{eq:rg-formation-fast-environment-coupling}
\end{equation}
For a parent at any $p\ge1$, couple its retained loop to the
auxiliary loop process. By
\cref{eq:rg-formation-parent-force-environment},
\begin{equation}
 |h_t-h_t^*|
 \le C_L\sqrt{TJ_1}
       \int_\sigma^t e^{-c_LJ_1(t-v)}d_\epsilon(y_v)dv.
 \label{eq:rg-formation-loop-environment-coupling}
\end{equation}
The environment in these two inequalities remains the actual
adapted environment. It has not been frozen or differentiated.
Multiplication by the row Lipschitz constant, with the direct
parent row error in
\cref{eq:rg-formation-parent-row-environment}, converts them to
proposal-row errors.

The bound needed at random proposal times is stronger than a
fixed-time unconditional moment. Condition on the whole independent
future Poisson grid, but not its marks or any future continuous
path. Brownian law is unchanged under this conditioning.
Evenness and the uniform lower curvature on $M_p$ give, on every
sheet and level,
\begin{equation}
 (T\Delta-\nabla V_{p,j}\cdot\nabla)|x|^2
 \le2Tk_0-2c_L|x|^2.
 \label{eq:rg-formation-grid-moment-generator}
\end{equation}
Proposals preserve $x$, and projections contract its norm.
The common actual preparation has second moment at most
$C_L(A^2+T)$. Thus
\begin{equation}
 \sup_{t\le H_*}\E[|X_{M_p,t}|^2\mid\text{Poisson grid}]
 \le C_L(A^2+T).
 \label{eq:rg-formation-grid-moment}
\end{equation}
If the common actual initial sample is also fixed, the right side
first contains its squared norm; integrating that sample returns
the displayed bound. The future grid after a deterministic
activation is independent of its actual past, so the same argument
applies at each activation.
In particular it applies at grid-selected times.

Average the row errors from
\cref{eq:rg-formation-fast-environment-coupling,eq:rg-formation-loop-environment-coupling}
using \cref{eq:rg-formation-grid-moment}. The exponential
convolution has integral at most $C_L/J_1$. Their common expected
per-proposal bound is
\begin{equation}
 \eta_\epsilon
 =C_L\left[\frac{J_2}{J_1}
       +\frac{J_2}{\sqrt{J_1}}
                       \left(1+\frac A{\sqrt T}\right)\right]
 =O_L(\epsilon^7).
 \label{eq:rg-formation-averaged-environment-error}
\end{equation}
Indeed $Q_1=2Q_2+20$ and $A/\sqrt T=\epsilon^{-3}$.
This averaging uses the actual distribution of every past prefix.
There is no division by a probability of a rare source prefix.

Define one comparison source process on the same proposal grid,
starting at the actual initial assignment and using row
$\mathcal B_{j\cdot}(\vartheta_0(t))$, where
$\vartheta_0(t)=2-e^{-\epsilon t}$.
This is a single comparison law common to both actual processes.
The actual dynamics continue to use their own laws and coefficients.
At a proposal time $t$, let $\sigma$ be the most recent preceding
proposal or activation. Condition on the common actual starting
sample, the full future grid, and all previous proposal marks,
including self marks. The assignment and active level are fixed on
$(\sigma,t)$; no future continuous trajectory has been conditioned
on. Starting the auxiliary diffusion from the actual fast point at
$\sigma$, use the environmental coupling, then
\cref{eq:rg-formation-fast-mixing-bound}, then the Lipschitz
comparison of $\vartheta(t)$ with $\vartheta_0(t)$.
For gaps at least
\begin{equation}
 h_*=\frac{k_L}{J_1}\log(1/\epsilon),
 \label{eq:rg-formation-history-gap}
\end{equation}
the averaged conditional row error is at most
\begin{equation}
 C_L\left[
 (1+A\sqrt{J_1/T})e^{-c_{\mathrm{fast},L}J_1h_*}
       +\epsilon/J_1+\eta_\epsilon+r_\epsilon\right].
 \label{eq:rg-formation-long-gap-row-error}
\end{equation}
The moment of the point-start mixing factor includes
$\sqrt{J_1/T}(A+\sqrt T)$; its $\sqrt{J_1}$ term is absorbed by
$A\sqrt{J_1/T}$ because $A/\sqrt T\ge1$.
After averaging over the actual past marks, their conditional
expectations return the grid-conditioned moments exactly.

Short gaps use the trivial total-variation bound one.
For a rate-one Poisson grid, the expected number of ordered pairs
of proposal times less than $h_*$ apart is at most $H_*h_*$,
by direct integration of the two-time intensity on
$\{0<s<t<H_*, t-s<h_*\}$. This bounds the count of short
consecutive gaps. Proposals within $h_*$ after the initial time or
one of the finitely many activations add at most $(L+1)h_*$.
The expected total proposal count is $H_*$.

For clarity, the passage from conditional row errors to the entire
history uses a finite-chain identity, rather than pointwise marginal
agreement. Conditional on a grid with $N$ proposals, write the two
joint mark laws as products of their conditional kernels.
Telescope by replacing kernels in order while keeping the earlier
prefix distributed according to the actual law and the later tail
according to the comparison law. Integration of the tail, whose
mass is one, gives
\begin{equation}
 \|\mathcal L(J_1,\ldots,J_N)
          -\mathcal L(\overline J_1,\ldots,\overline J_N)\|_{\mathrm{TV}}
 \le \sum_{m=1}^N
       \E_{\mathrm{actual\ prefix}}
       \|K_m(\cdot\mid J_{<m})
                 -\overline K_m(\cdot\mid J_{<m})\|_{\mathrm{TV}}.
 \label{eq:rg-formation-history-telescoping}
\end{equation}
The $J_m$ in this display are proposal marks, not energy strengths.
Integrating the grid and using the preceding bounds gives
\begin{align}
 \mathcal E_{\mathrm{hist}}\le C_L\big[
 &(H_*+L+1)h_*
 +H_*(1+A\sqrt{J_1/T})
                      e^{-c_{\mathrm{fast},L}J_1h_*}
 \notag\\
 &+H_*\epsilon/J_1
       +H_*\eta_\epsilon+H_*r_\epsilon\big].
 \label{eq:rg-formation-full-history-error}
\end{align}
With $c_{\mathrm{fast},L}k_L>Q_1/2+12$, the five contributions
are, respectively, bounded by constants times
$\epsilon^{19}\ell_\epsilon$, $\epsilon^8$,
$\epsilon^{Q_1}$, $\epsilon^6$, and
$\epsilon^{2Q_1+19}+\epsilon^{3Q_1+25}\ell_\epsilon^2$.
Thus $\mathcal E_{\mathrm{hist}}\le C_L\epsilon^6$.

For each common actual initial sample and grid, maximally couple
the two finite mark sequences to the comparison sequence.
Glue the two couplings through that common comparison law. This
preserves the actual source initial sample and its exact parent
projection. Fill in each continuous trajectory from its original
regular conditional distribution given its initial sample, grid,
and marks. Such distributions exist on finite-dimensional continuous
path spaces and finite-state right-continuous path spaces, which
are standard Borel spaces. Every actual full-process marginal is
therefore preserved. Forgetting self proposals makes the source
path a function of this grid and mark sequence and cannot increase
the disagreement probability. Summing over finitely many
activations changes only $C_L$.

This coupling need not be causal. It makes no total-variation claim
for the continuous trajectories. Continuous-path estimates can be
combined with it by retaining each original marginal path event;
they cannot be justified by treating Brownian motion conditioned on
future marks as a new independent Brownian motion.
\hfill$\square$

\section{Actual coherence, source movement, and continuous paths}
\label{sec:rg-formation-dynamics-transfer}

\lemmaheading{Successive coherence with the original conditional marks}
{lem:formation-coherence}
Under the preparation and regime of
\cref{eq:rg-formation-preparation-margin,eq:rg-formation-preparation-transport,eq:rg-formation-weakest-one-regime},
there are fixed positive constants $\eta_L,c_L,C_L$ and a finite
$k_{\mathrm{geom},L}$ such that the windows
$I_s=[t_s,2t_s]$, $t_s=k_L\log(1/\epsilon)/J_s$, with
$k_L>k_{\mathrm{geom},L}$, have the following property. Except on an
event of probability at most
$C_L(1+H_*)\exp(-c_L/\epsilon^6)$, the fine process and its actual
successive parent descriptions have coherent components
$3\cdot2^{L-1}\to2^{L-1}\to\cdots\to1$. On $I_s$, residuals of
levels at most $s$ and the local discrepancies $d_i$ are below the
fixed tolerance $\eta_LA$, whereas future bridges remain separated
by a fixed positive multiple of $A$. For eliminated coordinates this
statement uses the conditional root-mean-square original residuals
in \cref{eq:rg-formation-inherited-edge-statistic}, separately for
the physical posterior and restricted optimizer. Final coherence
persists through $H_*$. All loop coordinates remain retained. \status{ESTABLISHED}

\emph{Proof.}
Write the non-root Gaussian comparison of
\cref{lem:formation-prepared-continuation} on active level $p<s$
as a current normal component $u\in\mathcal N_p\cap\mathcal N_s^\perp$,
a future component $y\in\mathcal N_s$, and the separate loop
component. The weighted flag gives, for the current/future blocks
of $H_p$,
\begin{equation}
 F\succeq c_LJ_sI,
 \qquad \|G\|\le C_LJ_{s+1},
 \qquad \|D\|\le C_LJ_{s+1},
 \qquad J_{L+1}=0.
 \label{eq:rg-formation-current-future-blocks}
\end{equation}
Indeed rows of levels at most $s$ annihilate $\mathcal N_s$.
Consequently the original cross and future blocks contain only
strengths at most $J_{s+1}$. In an earlier elimination through $p$,
the eliminated-current block is $O_L(J_{p+1})$, its inverse hidden
block is $O_L(J_p^{-1})$, and the eliminated-future block is
$O_L(J_{s+1})$. Their Schur correction is therefore still
$O_L(J_{s+1})$, since $J_{p+1}\le J_p$. For $p=0$ there is no such
correction. Loop/normal cross blocks vanish in the Gaussian core.

Let $\mu^G$ be the deterministic reference mean. Its norm is at
most $C_LA$, and its loop component is zero. Integrating the current
equation
$\dot\mu_u^G=-\vartheta(t)(F\mu_u^G+G\mu_y^G)$, using
$1\le\vartheta\le21/10$, gives
\begin{equation}
 |\mu_u^G(t)|
 \le C_LA\left[e^{-c_LJ_s(t-t_0)}+\frac{J_{s+1}}{J_s}\right].
 \label{eq:rg-formation-current-mean}
\end{equation}
The full future row has norm $O_L(J_{s+1})$. Integrating that row
from the original common center and noting that earlier projections
preserve $P_s\mu^G$ gives, for a sequential process begun at
$\tau_{s-1}$,
\begin{equation}
 |\mu^G(t)-AP_sv_0|
 \le C_LA\left[
 e^{-c_LJ_s(t-\tau_{s-1})}
 +\frac{J_{s+1}}{J_s}+J_{s+1}t\right].
 \label{eq:rg-formation-window-mean}
\end{equation}
The future terms are absent when $s=L$. For a fine process use
$t_0=0$ in \cref{eq:rg-formation-current-mean}.

Here is the required uniform-in-time fluctuation bound. On each
active non-root space, $Z=X^G-\mu^G$ obeys
$dZ=-\vartheta(t)H_pZdt+\sqrt{2T}dW$, with
$\vartheta H_p\succeq h_LI$ and initial covariance at most
$C_LT H_p^{-1}$. Choose $a_*>0$ below $h_L$ and small enough for
the initial Gaussian exponential moment. For
$V(z)=\exp(a_*|z|^2/(4T))$, put $v=a_*|z|^2/(4T)$. Direct
differentiation gives
\begin{equation}
 \mathcal L_tV
 =\left[\frac{a_*d}{2}
 +\frac{a_*^2|z|^2}{4T}
 -\frac{a_*\vartheta z^TH_pz}{2T}\right]V
 \le a_*(d/2-v)e^v\le C_L.
 \label{eq:rg-formation-gaussian-exponential-generator}
\end{equation}
The initial expectation of $V$ is bounded. Orthogonal retention
decreases $V$, so the same estimate survives the finite sequence
of projections. Apply It\^o's formula up to the first exit from a
ball, localize the coefficients, and use the nonnegative stopped
value of $V$. Removing the localization yields, for every $\rho>0$,
\begin{equation}
 \Pr\left\{\sup_{t\le H_*}|X^G_t-\mu^G(t)|>\rho A\right\}
 \le C_L(1+H_*)
       \exp\left[-c_L\rho^2A^2/T\right].
 \label{eq:rg-formation-gaussian-maximum}
\end{equation}
Higher stiffnesses only strengthen the negative term in
\cref{eq:rg-formation-gaussian-exponential-generator}; no upper
stiffness enters this probability estimate.

Let $M_L$ bound the norms of the finitely many original residual
and $d_i$ covectors. Choose $\eta_L$ smaller than a fixed fraction
of $m_L$ and then a fixed fluctuation radius $\rho A$ so that
$M_L\rho A$ is smaller than the reserved tolerance margin.
Take, for example, $k_{\mathrm{geom},L}=2/c_L$ for the restoring
constant in \cref{eq:rg-formation-current-mean}; increasing this
finite constant is allowed. For sufficiently small $\epsilon$,
$\tau_{s-1}<t_s/2$, so the exponential term in
\cref{eq:rg-formation-window-mean} tends to zero uniformly on
$I_s$. Also
\begin{equation}
 \frac{J_{s+1}}{J_s}\longrightarrow0,
 \qquad J_{s+1}t_s\longrightarrow0\quad(s<L),
 \qquad \frac{\tau_L}{H_*}\longrightarrow0.
 \label{eq:rg-formation-geometric-entrance}
\end{equation}
The nonlinear correction from
\cref{eq:rg-formation-path-force-bound} is at most
$C_LA\epsilon^{5/2}$, including the allowed initial transport and
finitely many inherited errors. This bound holds for every adapted
assignment path. Since $A^2/T=\epsilon^{-6}$, the strict prepared
residual margins and \cref{eq:rg-formation-gaussian-maximum} now
give the asserted probability and separation on every window.
The stable current-block equation continues to bound previously
formed residuals as their future forcing decreases. After the last
normal block disappears there is no future forcing. Thus the final
coherent regime persists through $H_*$ on the same event.

It remains to check what an eliminated edge means. For a cumulative
elimination through $p\ge1$, write a fine non-root state as $z+y$,
where $z\in W_p^\perp$ and $y\in M_p$. The optimizing Gaussian
precision is $\vartheta H$; the physical one is
$\vartheta H-TR$. Every row of $b_d$ belongs to the first normal
layer, so $R=b_d^Tb_d/(1-c)$ annihilates $\mathcal N_1$,
$\mathcal H$, and $\mathcal K$. The correction therefore has no
hidden/future or hidden/loop block. The admitted inequality
$TR\preceq H/2$ preserves a hidden precision lower bound
$c_LJ_pI$ for both laws. Their hidden/future block is
$O_L(J_{p+1})$, while the hidden/loop block is zero.

Consequently the corresponding Gaussian conditional mean has norm
at most $C_L(J_{p+1}/J_p)|y_{\mathcal N_p}|$ and its covariance
trace is at most $C_LT/J_p$. Conditional on the full assignment,
each exact hidden density is a normalized tilt of its own Gaussian
conditional by a factor with log oscillation at most $B_0$.
Its normalized density ratio is at most $e^{B_0}$. For every
retained state with $|y_{\mathcal N_p}|\le C_LA$,
\begin{equation}
 \E_{\mu_\vartheta}
   [|z|^2\mid P_pX=y,S=j]
 \le C_Le^{B_0}\left[
       A^2\left(\frac{J_{p+1}}{J_p}\right)^2
       +\frac{T}{J_p}\right],
 \quad
 \mu_\vartheta\in
 \{\boldsymbol\Pi_{\vartheta,\mathsf o},\pi_\vartheta\}.
 \label{eq:rg-formation-hidden-original-mark-moment}
\end{equation}
The normal-state bound holds on the event already constructed.
For every original residual covector $r_e$, the triangle inequality
and its reverse in conditional $L^2$ give
\begin{equation}
 \left|
 \sqrt{\E_{\mu_\vartheta}[r_e(X)^2\mid P_pX=y,S=j]}
 -|r_e(y)|\right|
 \le\|r_e\|
      \sqrt{\E_{\mu_\vartheta}[|z|^2\mid P_pX=y,S=j]}.
 \label{eq:rg-formation-original-mark-rms}
\end{equation}
The right side is $o(A)$ under the common regime, uniformly in
the full assignment and separately for the two laws. The same
argument applies to $d_i$. It transfers both the small formed
residuals and the separated future residuals to
\cref{eq:rg-formation-inherited-edge-statistic}. No erased mark is
replaced by its value at $P_px$, and no complete mixture KL is
identified with a scalar residual. The bridge levels join successive
pairs of already formed components, giving the stated counts.
\hfill$\square$

\lemmaheading{Actual source envelope and changing hidden kernels}
{lem:formation-source-envelope}
Every actual process or finite replacement lineage in
\cref{lem:formation-prepared-continuation} satisfies
\cref{eq:rg-formation-source-movement} for the distinguished binary
source, measured from any of its actual activation times. Between
$\vartheta=1$ and its final learned parameter, the first hidden
physical conditional and the first hidden optimizing conditional
each change by KL greater than $1/200$, uniformly at the same
retained input and full assignment. \status{ESTABLISHED}

\emph{Proof.}
At the distinguished receiver the posterior log odds equal
\begin{equation}
 \log\frac{\beta^P_{i_0,+}(x,\vartheta)}
               {\beta^P_{i_0,-}(x,\vartheta)}
 =4(\vartheta-3/2)
  -\delta\left[
 f\bigl(\sqrt{J_1/T}r_{i_0,+}(x)\bigr)
 -f\bigl(\sqrt{J_1/T}r_{i_0,-}(x)\bigr)\right].
 \label{eq:rg-formation-source-log-odds}
\end{equation}
Since $0\le f\le1$, this row lies between
$\sigma(4(\vartheta-3/2)-\delta)$ and
$\sigma(4(\vartheta-3/2)+\delta)$, where
$\sigma(u)=(1+e^{-u})^{-1}$. Summing the exact inherited proposal
row over all assignments with distinguished label $+$ and taking
the actual expectation preserves these bounds. Apply the generator
to the indicator of that label. With
$\lambda(t)=\Pr\{S_{i_0}(t)=+\}$, the result is
\begin{equation}
 \lambda'=b_\lambda-\lambda,
 \qquad
 \sigma(4(\vartheta-3/2)-\delta)
 \le b_\lambda(t)\le
 \sigma(4(\vartheta-3/2)+\delta).
 \label{eq:rg-formation-source-forcing-envelope}
\end{equation}
Neither the diffusion nor a marked projection changes this
indicator. Thus variation of constants remains valid across every
activation, even when the inherited row depends on the current
assignment and continuous state.

By \cref{lem:formation-parameter-comparison}, uniformly before
$\tau_L=2k_L\log(1/\epsilon)$,
$\vartheta\le1+C_L\epsilon\log(1/\epsilon)+C_Lr_\epsilon$.
A solution of \cref{eq:rg-formation-source-forcing-envelope} cannot
exceed the larger of its starting value and the supremum forcing.
The preparation tolerance and $\|\sigma'\|_\infty\le1/4$ give
\begin{equation}
 \lambda(\tau_s)
 \le\sigma(-2)+10^{-2}+\delta/4
       +C_L[\epsilon\log(1/\epsilon)+r_\epsilon],
 \qquad 0\le s\le L.
 \label{eq:rg-formation-source-at-activation}
\end{equation}
There is no estimate on the number of proposals in this argument.
For small $\epsilon$, every activation precedes $H_*-10$ and
$\vartheta\ge7/4$ throughout $[H_*-10,H_*]$. Nonnegativity of
the earlier contribution in variation of constants yields
\begin{equation}
 \lambda(H_*)\ge(1-e^{-10})\sigma(1-\delta).
 \label{eq:rg-formation-source-final-mass}
\end{equation}
For $\delta<10^{-4}$, the elementary estimates
$e^{1-\delta}>249/100$, $e^2>7$, and $e^{10}>2^{10}$ imply
\begin{equation}
 (1-e^{-10})\sigma(1-\delta)-\sigma(-2)
       -10^{-2}-\delta/4
 >\frac{1023}{1024}\frac{249}{349}
       -\frac18-\frac1{100}-\frac1{40000}
 >\frac{57}{100}.
 \label{eq:rg-formation-source-strict-margin}
\end{equation}
The remaining error in
\cref{eq:rg-formation-source-at-activation} tends to zero. Hence
$\lambda(H_*)-\lambda(\tau_s)>1/2$ for every admitted activation.
For a binary row this is its total-variation change. The calculation
uses the actual recognition marginal throughout.

For the hidden-kernel statement, fix the retained input and full
assignment, and first omit the bounded source factor. On the first
hidden normal space, let $D_1,D_f$ be the initial and final Gaussian
precisions, with covariances $TD_1^{-1},TD_f^{-1}$. For the optimizer
the precision scales with $\vartheta$. For the physical law,
on the non-root space,
\begin{equation}
 D_1=H-TR\succ0,
 \qquad
 D_f=D_1+(\vartheta_f-1)H
       \succeq\vartheta_fD_1,
 \qquad \vartheta_f\ge7/4.
 \label{eq:rg-formation-physical-precision-growth}
\end{equation}
Restriction to a conditional block preserves the order. If an
earlier Gaussian marginal is taken, its Schur quadratic is the
minimum over the discarded coordinates, which also preserves this
order. In either case the weaker common bound
$D_f\succeq(5/4)D_1$ holds. The conditional Gaussian means may
differ; equality of means is not needed.

For Gaussian laws $G_1,G_f$ on a hidden space of dimension $k\ge1$,
whiten $D_1$ and diagonalize the conjugated final precision, whose
eigenvalues satisfy $\rho_i\ge5/4$. Completing the Gaussian square
in their Hellinger affinity gives
\begin{equation}
 \mathcal A(G_1,G_f)
 =\prod_{i=1}^k
       \left(\frac{2\sqrt{\rho_i}}{1+\rho_i}\right)^{1/2}
       e^{-M},
 \qquad M\ge0.
 \label{eq:rg-formation-gaussian-affinity}
\end{equation}
The factor decreases for $\rho_i\ge1$, so this is at most
$[2\sqrt{5/4}/(1+5/4)]^{k/2}$. For each of the two exact
conditionals $k_1,k_f$, its normalized source tilt relative to its
own Gaussian conditional is bounded above by $e^{B_0}$, including
after exact hidden integration. Therefore
$\mathcal A(k_1,k_f)\le e^{B_0}\mathcal A(G_1,G_f)$.
Jensen's inequality under $k_1$ gives
$-\KL(k_1\Vert k_f)/2\le\log\mathcal A(k_1,k_f)$ and consequently
\begin{equation}
 \KL(k_1\Vert k_f)
 \ge\frac{k}{2}\log\frac{81}{80}-2B_0
 \ge\frac1{162}-\frac1{5000}>\frac1{200}.
 \label{eq:rg-formation-hidden-kernel-change}
\end{equation}
Here $\log(1+1/80)\ge1/81$. The means and bounded tilts make the
likelihood-ratio use integrable. This argument applies separately
to the physical and optimizing conditionals, uniformly in retained
input and assignment. It does not reset an actual hidden law.
\hfill$\square$

\lemmaheading{Continuous displacement compatible with source histories}
{lem:formation-continuous-path}
For an actual source and its exact parent started from a common
actual retained sample and parameter at $\tau_s$, the coupling of
\cref{lem:formation-history-coupling} can be extended without
changing either full process marginal so that
\begin{equation}
 \Pr\left\{\sup_{\tau_s\le t\le H_*}
 |P_sX_t-\overline X_t|>C_LA\epsilon^2\right\}
 \le C_L(1+H_*)e^{-c_L/\epsilon^2}.
 \label{eq:rg-formation-continuous-displacement-proof}
\end{equation}
Its common actual activation information and complete assignment
history bound are retained. A finite number of successive
replacements changes only the constants. \status{ESTABLISHED}

\emph{Proof.}
First compare Gaussian means at the same deterministic parameter
path. In the current/future split of
\cref{eq:rg-formation-current-future-blocks}, the inherited normal
precision is $H_s=D-G^TF^{-1}G$. Put
$v=\mu_u^G+F^{-1}G\mu_y^G$. The source future mean satisfies
\begin{equation}
 \dot\mu_y^G=-\vartheta H_s\mu_y^G-\vartheta G^Tv,
 \qquad
 |v(t)|\le C_LA\left[
 e^{-c_LJ_s(t-\tau_s)}+\frac{J_{s+1}}{J_s}\right].
 \label{eq:rg-formation-gaussian-mean-defect}
\end{equation}
The second bound follows from
\cref{eq:rg-formation-current-mean} and
$\|F^{-1}G\|\le C_LJ_{s+1}/J_s$. Integrating the contractive
future propagator from a matching retained mean gives
\begin{equation}
 \sup_{\tau_s\le t\le H_*}
 |\mu_y^G(t)-\overline\mu^G(t)|
 \le C_LA\left[
 \frac{J_{s+1}}{J_s}
 +H_*\frac{J_{s+1}^2}{J_s}\right].
 \label{eq:rg-formation-gaussian-mean-reduction}
\end{equation}
There is no future normal contribution when $s=L$. The loop
Gaussian core is unchanged by the reduction and has zero mean.

Now allow the independently learned parameters to differ. Write
$\chi=\sup|\vartheta-\overline\vartheta|\le C_Lr_\epsilon$ by
\cref{lem:formation-parameter-comparison}. For homogeneous parent
means $m,\overline m$ and their difference $d$, Young's inequality
in the $H_s$ quadratic form gives
\begin{equation}
 \frac d{dt}|d|^2
 \le-\vartheta d^TH_sd
       +C_L\chi^2\overline\vartheta
                         \overline m^TH_s\overline m.
 \label{eq:rg-formation-mean-parameter-energy}
\end{equation}
Since
$\frac d{dt}|\overline m|^2
=-2\overline\vartheta\overline m^TH_s\overline m$,
the time integral of the final energy is at most $C_LA^2$.
The extra displacement is therefore at most $C_LA\chi$; no
largest-stiffness factor times the horizon is introduced.

Attach to each correct actual marginal its own Gaussian comparison
from \cref{lem:formation-prepared-continuation}. At activation its
reference starts from the projected preceding Gaussian reference,
while the actual parent starts from the full actual retained
sample. Their inherited discrepancy is already included in
\cref{eq:rg-formation-path-force-bound}. Apply
\cref{eq:rg-formation-gaussian-maximum} to each centered reference
with $\rho=\epsilon^2$, and then use the two nonlinear force
bounds. The triangle inequality gives, outside a set of probability
$C_L(1+H_*)e^{-c_L/\epsilon^2}$,
\begin{equation}
 \sup_{\tau_s\le t\le H_*}|P_sX_t-\overline X_t|
 \le C_LA\left[
 \epsilon^2+\epsilon^{5/2}+r_\epsilon
 +\frac{J_{s+1}}{J_s}
 +H_*\frac{J_{s+1}^2}{J_s}\right].
 \label{eq:rg-formation-continuous-error-terms}
\end{equation}
For $s<L$, the last term equals $2\epsilon^{19}$ and
$J_{s+1}/J_s=\epsilon^{Q_{s+1}+20}$. The quantity
$r_\epsilon=A^2J_1+T[1+\log(1/\epsilon)]^2$ is $o(\epsilon^2)$.
Both future terms vanish at $s=L$. This proves
\cref{eq:rg-formation-continuous-displacement-proof} for the
non-root coordinates. The autonomous root law and dynamics are
identical and independent of the source history, so their root
coordinates can be coupled identically.

Finally take the coupling of complete source histories and common
activation information from \cref{lem:formation-history-coupling}.
Extend it by the conditional law of each actual path together with
its own reference, given that process's history and activation
information. These conditional probability kernels exist on the
finite-dimensional continuous-path and finite-state path spaces,
which are standard Borel. Each process-reference pair retains its
correct marginal joint law. The force comparison holds for every
adapted assignment path, and each reference concentration event
retains its marginal probability under this extension. Their union
bound therefore gives the same displacement estimate on the
history coupling. This argument does not declare Brownian motions
to remain shared after conditioning on source histories. Common
activation samples are preserved, and no claim of total variation
convergence of continuous path laws is required. A finite sequence
of replacements uses a finite triangle inequality and union bound.
\hfill$\square$

\section{Joining the estimates in one regime}
\label{sec:rg-formation-common-regime-proof}

The proofs use the single regime
$Q_L=0$, $Q_s=2Q_{s+1}+20$, $J_s=\epsilon^{-Q_s}$,
$A^2=\epsilon^{3Q_1+20}$, and $T=A^2\epsilon^6$.
Choose the fixed entrance constant to satisfy
\begin{equation}
 k_L>
 \max\left\{
 k_{\mathrm{geom},L},
 \frac{2(Q_1+20)}{c_{\mathrm{early},L}},
 \frac{Q_1+8}{c_{\mathrm{joint},L}},
 \frac{Q_1/2+12}{c_{\mathrm{fast},L}}
 \right\}.
 \label{eq:rg-formation-joined-entrance-proof}
\end{equation}
Each denominator is positive for this finite graph and fixed
nonlinear coefficient. The first condition supplies the geometric
windows, the second the early conditional-entropy entrance, the
third terminal joint relaxation, and the fourth the complete
source-history estimate. For sufficiently small $\epsilon$ these
windows are ordered, their activations precede $H_*$, and the same
admitted preparation satisfies every estimate. Thus the separate
bounds prove the conjunction in
\cref{thm:finite-depth-adaptive-formation}; no estimate uses a
different temperature, source rate, or initialization.
\status{ESTABLISHED}

Normalization, marked pushforward composition, inherited proposal
rows, and the objective chain rule are exact identities at each
positive parameter instance. Formation probabilities, conditional
entropy, complete covector error, parameter comparison, and path
couplings are asymptotic bounds in the displayed regime. The
positive nonlinear coefficient and source proposal rate remain
fixed, as does $J_L=1$. All constants and neighborhoods may depend
on the finite depth. Extension to mixed-level source support,
uniform depth, or base-manifold transport is not proved here.
\status{NOT-CLAIMED}
